\section{Result, comparison target, and limits}

\paragraph{The result.}
General numerical equality transport can be derived from a rule that
substitutes only for a \emph{free numerical variable} inside a term, together
with equality composition at the root. The argument gives a fresh local
name to the term to be transported. The logical calculus itself supplies
that name and later discharges it. Consequently the trusted equality
instruction receives two proofs and no term template. It does not search
for occurrences of a compound source term.

The resulting theory retains raw iterator syntax and ordinary
capture-avoiding substitution. Its only public formula constructors are
equality and the paired binder
\[
 D_{x,y}B := \exists x\bigl(B(x,0)\to\forall y\,B(x,y)\bigr).
\]
There are five serialized instruction families: Compute, Param, Use,
Open, and Close. Compute contains the two usual iterator equations;
there is no primitive reflexivity instruction. The detailed statement
below is the complete axiom-and-rule specification, including every
side condition and each branch of the finite operations.

\paragraph{The equivalence established.}
The paper proves exact preservation and reflection of ordinary sequents,
and both formula round trips, between this presentation and the native
calculus in the supplied repository. Composing this comparison with its
formalized arithmetic bridge establishes bi-interpretability with full
first-order intuitionistic PRA. Both numerical domains are whole, equality
is actual equality, and the numerical composite isomorphisms have graph
$x=y$. This is stronger than agreement on true numerical equations in
the standard natural numbers.

\paragraph{Pinned baseline.}
The source snapshot is ProveIt commit
\texttt{5d85023b0688e03b75f0ec3a04bb43ef0bb51705}, under
\texttt{Logic/PrimitiveRecursiveArithmetic/ConciseIPRA}.
The review included its index of 56 filed packages and the five archives
then present under \texttt{docs/incoming}. The relevant predecessors are
the raw clause presentation, template-free equality, one-term transport,
and the incoming root-free and uniform-step presentations
\cite{index,rawclausal,templatefree,incoming}.

The clause interface and fresh-alias argument are inherited ingredients.
The specific reduction isolated here is that the fresh-alias argument
requires proper-subterm replacement only at variables: replacement of
compound patterns can be removed as well. The written proof and the
executable proof compiler make that restriction explicit. No claim of
priority outside the inspected repository is made.

\paragraph{What ``simpler'' means here.}
The equality interface is mathematically weaker and has no supplied
template. In a matched implementation with unused machinery removed
from the control, the complete program decreases from 2907 to 2857
normalized bytes. It uses four more significant tokens and eleven more
AST nodes. Thus the result is a structural reduction and a modest byte
improvement, not a simultaneous improvement in every metric.
The original filed raw clause checker is 3046 bytes; comparisons with
it also include implementation changes and must not all be credited to
the equality theorem.

There is no established absolute minimum. The repository's compact
codecs remain much smaller. Moreover, an implementation refinement of
the existing two-case Eval calculus measures 2731 bytes in this same
investigation. The new presentation is therefore recommended for its
restricted equality rule and ordinary substitution, not advertised as
the smallest known complete Python program.

\paragraph{Proof status.}
All new mathematical implications used in the construction are written
out below. The complete Python checker and proof transformations have
been executed. The new kernel has not been formalized in Lean or Rocq.
The inherited native/iPRA equivalence is a separately formalized
dependency, whose relevant source modules and theorem names are included.

